\paragrafos{Προβλήματα αναλογιών: από το ένα στο πολλά}{Α.6.5, «Προβλήματα αναλογιών»}

\begin{anagke}
Η Α3 ετοιμάζει ένα μεγάλο φαγητό για το παζάρι. Η συνταγή της γιαγιάς Μυρτούς γράφει «για $4$ άτομα: $300$ γραμμάρια ρύζι». Θα έρθουν $10$ άτομα. «Πόσο ρύζι θέλουμε;» ρωτάει η Ρήνη. Ο Βασίλης προτείνει: «Για τέσσερα $300$, για άλλα τέσσερα $300$, και για τα υπόλοιπα δύο κάτι λιγότερο. Να βάλουμε $800$». Η Πολυξένη τον σταματάει: «Μην μαντεύεις. Πόσο ρύζι τρώει ένα άτομο;». Η Φανή λέει ότι ένα άτομο θέλει $300:4=75$ γραμμάρια, άρα τα $10$ άτομα $750$. «Αυτό δεν είναι ο κανόνας των τριών που έλεγαν στο δημοτικό;» ρωτάει ο Στέλιος. Η Αγγελική παρατηρεί κάτι ακόμη: «Δεν είναι όλα τα προβλήματα έτσι. Αν η Ζωή είναι $5$ χρονών όταν η αδελφή της είναι $2$, δεν σημαίνει ότι όταν η Ζωή γίνει $10$ η αδελφή της θα γίνει $4$». Η παρέα καταλαβαίνει ότι πριν από οποιαδήποτε μέθοδο πρέπει να ελέγξει αν τα ποσά είναι ανάλογα.

Σε αυτή την παράγραφο η παρέα μαζεύει όλες τις τεχνικές που ξέρει: πίνακα, αναλογία, αναγωγή στη μονάδα, γράφημα, και χρησιμοποιεί και τον έλεγχο της αναλογίας ώστε να μην εφαρμόζει κανόνες όπου δεν ισχύουν.
\end{anagke}

\begin{lexi}{αναγωγή στη μονάδα}{αναγωγή στη μονάδα}
\lexiword{Αναγωγή} σημαίνει «φέρνω πίσω»: φέρνω το πρόβλημα πίσω στη μονάδα, δηλαδή βρίσκω πρώτα πόσο αντιστοιχεί στο ένα, και μετά πολλαπλασιάζω. Το \lexiword{ανά} σημαίνει «για κάθε»: «ανά άτομο» είναι «για κάθε άτομο».
\end{lexi}

\begin{metafrasi}
Για $4$ άτομα $300$ g, για $10$ άτομα $x$ g. & $\dfrac{4}{10}=\dfrac{300}{x}$\\
Πόσο για ένα άτομο; & $300:4=75$ g\\
Για $10$ άτομα. & $10\cdot75=750$ g\\
Ο πίνακας των δύο ποσών. & $\begin{array}{c|cc}\text{άτομα}&4&10\\\hline\text{ρύζι}&300&x\end{array}$\\
\end{metafrasi}

\tmima{Η θεωρία}

\tmimasmall{Πρώτα ελέγχω αν είναι ανάλογα}

Για να δω αν δύο ποσά είναι ανάλογα χρησιμοποιώ έναν από τους τρεις τρόπους της προηγούμενης παραγράφου: τον ορισμό (πολλαπλασιάζω το ένα, πολλαπλασιάζεται και το άλλο), τον τύπο $y=\alpha x$ ή τον σταθερό λόγο $\frac{y}{x}$. Στα προβλήματα της ζωής σκέφτομαι: \emph{αν διπλασιάσω το ένα, διπλασιάζεται και το άλλο;} Αν η απάντηση είναι ναι, τα ποσά είναι ανάλογα.

\tmimasmall{Τρεις τρόποι επίλυσης}

\begin{enumerate}[leftmargin=1.6em,itemsep=0.3ex,topsep=2pt]
\item \textbf{Αναγωγή στη μονάδα.} Βρίσκω πρώτα την τιμή για το $1$. $300:4=75$ g το άτομο, και $10\cdot75=750$ g.
\item \textbf{Αναλογία (τέταρτος ανάλογος).} Οργανώνω πίνακα και γράφω ισότητα λόγων: $\dfrac{4}{10}=\dfrac{300}{x}$, άρα $4x=3\,000$, $x=750$.
\item \textbf{Συντελεστής αναλογίας.} Βρίσκω τον συντελεστή $\alpha=\frac{300}{4}=75$ και γράφω $y=75x$. Για $x=10$: $y=750$. Ή από τη γραφική παράσταση.
\end{enumerate}

Και οι τρεις δίνουν το ίδιο. Διαλέγω αυτόν που μου φαίνεται πιο εύκολος και, αν μπορώ, ελέγχω με έναν δεύτερο.

\tmimasmall{Διαιρώ ένα ποσό ανάλογα με δοσμένους αριθμούς}

Όταν μοιράζω ένα ποσό ανάλογα με τους αριθμούς $3$ και $5$, σημαίνει ότι το κάθε μέρος είναι $3$ και $5$ «μονάδες» αντίστοιχα. Αν το σύνολο είναι $120$, οι μονάδες είναι $3+5=8$, η μία μονάδα $120:8=15$ και τα μέρη $3\cdot15=45$ και $5\cdot15=75$.

\begin{prosoxi}
\begin{itemize}[leftmargin=1.2em,itemsep=0.3ex,topsep=2pt]
\item Δεν είναι όλα ανάλογα. Αν η Ζωή είναι $5$ χρονών και η αδελφή της $2$, οι ηλικίες διαφέρουν πάντα κατά $3$: όχι ανάλογα.
\item Ο ίδιος αριθμός ($4$) δεν είναι πάντα το $x$: διάβασε προσεκτικά ποιο ποσό ποιον αριθμό δίνει.
\item Στον πίνακα βάζω τα ομοειδή μεγέθη στην ίδια γραμμή: άτομα με άτομα, γραμμάρια με γραμμάρια.
\item Αν βγει αποτέλεσμα που δεν έχει νόημα (π.χ. $3{,}5$ άνθρωποι), το κρατάω ως τιμή αλλά αναθεωρώ την απάντηση για τη ζωή.
\end{itemize}
\end{prosoxi}

\tmima{Λυμένα παραδείγματα}

\paradeigma{Το ρύζι της συνταγής}
\tixero Για $4$ άτομα η συνταγή θέλει $300$ g ρύζι και $2$ κρεμμύδια.\par
\tithelo Η ποσότητα για $10$ άτομα.\par
\possindeo Αναγωγή στη μονάδα: ρύζι ανά άτομο $300:4=75$ g, άρα $750$ g. Κρεμμύδια ανά άτομο $2:4=0{,}5$, άρα $5$.\par
\apantisi $750$ g ρύζι και $5$ κρεμμύδια.\par
\elegxos Με αναλογία: $\frac{4}{10}=\frac{300}{x}$ δίνει $x=750$.

\paradeigma{Το ύφασμα}
\tixero Τα $5$ μ. ύφασμα κοστίζουν $36$ €.\par
\tithelo Πόσο κοστίζουν $8$ μ.\par
\possindeo Ένα μέτρο κοστίζει $36:5=7{,}2$ €. Τα $8$ μ. κοστίζουν $8\cdot7{,}2=57{,}6$ €.\par
\apantisi $57{,}60$ €.\par
\elegxos Με αναλογία $\frac{5}{8}=\frac{36}{x}$, $5x=288$, $x=57{,}6$.

\paradeigma{Το ταξίδι}
\tixero Ένα αυτοκίνητο με σταθερή ταχύτητα κάνει $120$ km σε $1{,}5$ ώρες.\par
\tithelo Σε πόση ώρα θα κάνει $280$ km.\par
\possindeo Η απόσταση και ο χρόνος είναι ανάλογα. Ταχύτητα $120:1{,}5=80$ km/h. Χρόνος $280:80=3{,}5$ ώρες.\par
\apantisi $3{,}5$ ώρες.\par
\elegxos $80\cdot3{,}5=280$.

\paradeigma{Δεν είναι ανάλογα}
\tixero Η Ζωή είναι $5$ ετών όταν η αδελφή της είναι $2$.\par
\tithelo Πόσο θα είναι η αδελφή όταν η Ζωή είναι $15$;\par
\possindeo Αν ήταν ανάλογα θα είχαμε $6$. Αλλά οι ηλικίες διαφέρουν πάντα κατά $3$ χρόνια: $15-3=12$.\par
\apantisi $12$ ετών.\par
\elegxos Οι λόγοι: $\frac{2}{5}=0{,}4$ και $\frac{12}{15}=0{,}8$. Δεν είναι ίσοι.

\paradeigma{Ανάλογη μοιρασιά}
\tixero Ο Νίκος δούλεψε $3$ ώρες και ο Βασίλης $5$ ώρες στο παζάρι και μοιράζονται $120$ € ανάλογα με τις ώρες.\par
\tithelo Πόσα παίρνει ο καθένας.\par
\possindeo Οι ώρες είναι $3+5=8$. Η κάθε ώρα αντιστοιχεί σε $120:8=15$ €. Νίκος $3\cdot15=45$ €, Βασίλης $5\cdot15=75$ €.\par
\apantisi $45$ € και $75$ €.\par
\elegxos $45+75=120$ και $\frac{45}{3}=\frac{75}{5}=15$.

\tmima{Ασκήσεις}

\zoni{Το κατάλαβα}{αναγωγή στη μονάδα και αναλογία}

\ex Τα $3$ μπουκάλια νερό κοστίζουν $2{,}4$ €. Πόσο κοστίζουν τα $8$;

\ex Τα $5$ κιλά πορτοκάλια κοστίζουν $6{,}5$ €. Πόσο κοστίζουν $8$ κιλά; Πόσα κιλά παίρνεις με $20$ €;

\ex Τα $7$ τεμάχια ενός προϊόντος κοστίζουν $31{,}5$ €. Πόσο κοστίζουν τα $12$;

\ex Συμπλήρωσε τον πίνακα αναλόγων ποσών: $x$: $3,\ 5,\ 8,\ 12$ και $y$: $9,\ ?,\ ?,\ ?$.

\ex Ποια ποσά είναι ανάλογα; α) ηλικία και ύψος \quad β) κιλά και κόστος \quad γ) πλευρά και περίμετρος ισόπλευρου τριγώνου \quad δ) ώρες εργασίας και αμοιβή με σταθερό μεροκάματο \quad ε) πλευρά και εμβαδόν τετραγώνου \quad στ) ηλικία πατέρα και ηλικία γιου.

\zoni{Το χρησιμοποιώ}{κουζίνα, ταξίδια, δουλειές}

\ex \textbf{Το ρύζι.} Για $4$ άτομα η συνταγή θέλει $300$ g ρύζι και $2$ κρεμμύδια. Βρες τις ποσότητες για $10$ άτομα και για $6$ άτομα.

\ex \textbf{Οι πίτες.} Η γιαγιά Μυρτώ φτιάχνει $6$ πίτες από $1{,}2$ κιλά αλεύρι. Πόσο αλεύρι θέλει για $15$ πίτες;

\ex \textbf{Το ταξίδι.} Ένα αυτοκίνητο κάνει $120$ km σε $1{,}5$ ώρες. Πόσα km θα κάνει σε $3{,}5$ ώρες με την ίδια ταχύτητα;

\ex \textbf{Η κατανάλωση.} Ένα αυτοκίνητο καταναλώνει $8$ λίτρα ανά $100$ km. Πόση βενζίνη θέλει για $340$ km;

\ex \textbf{Οι φωτοτυπίες.} Οι $35$ φωτοτυπίες κοστίζουν $2{,}45$ €. Πόσο κοστίζουν οι $120$;

\ex \textbf{Οι ελιές.} Τα $15$ ελαιόδεντρα δίνουν $240$ κιλά ελιές. Πόσα κιλά δίνουν $40$ δέντρα;

\ex \textbf{Το νόμισμα.} Με ισοτιμία $1$ € $=1{,}1$ δολάρια, πόσα δολάρια είναι τα $50$ €; Πόσα ευρώ είναι τα $132$ δολάρια;

\ex \textbf{Η μοιρασιά.} Ο Νίκος δούλεψε $3$ ώρες και ο Βασίλης $5$ ώρες στο παζάρι, και μοιράζονται $120$ € ανάλογα με τις ώρες. Πόσα παίρνει ο καθένας;

\zoni{Το σκέφτομαι}{όταν ο κανόνας δεν ισχύει}

\ex Η Φανή είναι $5$ ετών όταν η αδελφή της είναι $2$. Όταν η Φανή είναι $15$, πόσο είναι η αδελφή της; Γιατί δεν είναι ανάλογα;

\ex Ο Ηλίας τρέχει $5$ km σε $25$ λεπτά. Θα τρέξει $10$ km σε $50$ λεπτά; Συζήτησε πότε το συμπέρασμα ισχύει και πότε όχι.

\ex Λύσε με δύο τρόπους: τα $9$ κιλά κοστίζουν $22{,}5$ €. Πόσο κοστίζουν $14$ κιλά;

\ex Ο Στέλιος λέει: «$3$ εργάτες χτίζουν έναν τοίχο σε $6$ ημέρες, άρα $6$ εργάτες σε $12$ ημέρες». Έχει δίκιο; Πόσες ημέρες περιμένεις να χρειαστούν $6$ εργάτες;

\ex Φτιάξε ένα δικό σου πρόβλημα αναλογίας με τιμή ανά μέτρο και λύσε το με τρεις τρόπους.

\zoni{Για το σπίτι}{στην κουζίνα}

\ex Πάρε μια συνταγή από το σπίτι σου και μετέτρεψέ την για διπλάσιο αριθμό ατόμων και για μισό αριθμό ατόμων.

\ex Βρες την τιμή ενός προϊόντος ανά κιλό και ανά $100$ γραμμάρια και σύγκρινε δύο συσκευασίες.
